\setlength\cwii{\b}
\setlength\cwiii{\c}
\setlength\cwi{\a}
\begin{tcbitemize}[raster columns=3,raster equal height=rows,
        raster force size=false,
        raster equal skip=-1.0pt,
        raster column 1/.style={width=\cwi,
            fontupper=\bfseries},
        raster column 2/.style={width=\cwii},
        raster column 3/.style={width=\cwiii,fontupper=\bfseries,
            halign=center,valign=center},
        ]
    %----------------------------------------------------------------
    \tcbitem 1)
    \tcbitem 
        \begin{figure}[H]
            \centering
            \begin{circuitikz}
                \newdimen\d \d=3cm
                \draw (0,0) coordinate (A)
                    node[left] {$\mathrm{A}$}
                    to[vsource,v=$E$,*-] ++(0,\d)
                    to[cute closed switch,l_=$\mathrm{K}$,-*] ++(\d,0)
                    coordinate (D)
                    node[above] {$\mathrm{D}$}
                    to[R,l_=$R$,v^<=$u_{R}$,-*] ++(0,-\d) coordinate (B)
                    node[below] {$\mathrm{B}$}
                    to[C,l_=$C$,v^<=$u_{C}$,i>_=$i$] (A);
                \node[eground] at (A) {};
            \end{circuitikz}
        \end{figure}
    \tcbitem (1pt)
    %----------------------------------------------------------------
    \tcbitem[raster multirow=4] 2)
    \tcbitem[raster multicolumn=2,raster multirow=2,blankest]
        \begin{tcbitemize}[raster rows=2,raster columns=2,
               raster height=\tcbtextheight,
               raster column 1/.style={width=\cwii},
               raster column 2/.style={width=\cwiii,fontupper=\bfseries}
            ]
            \tcbitem $u_{C}$ correspond à la courbe~\ding{182}.
            \tcbitem (0,5pt)
            %--------------------------------------------------------
            \tcbitem \textbf{Justification~:} 
                     À ${t=0}$, le condensateur est déchargé \
                     $\Rightarrow u_{C}(0)=0$.
            \tcbitem (0,5pt)
        \end{tcbitemize}
    %----------------------------------------------------------------
    \tcbitem 3)
    \tcbitem Loi des mailles~: \ $u_{C}+u_{R}=E$
             $\quad\Rightarrow\quad$ $u_{C}+R\times i=E$ \ (loi d'Ohm)
    
             Or $i=C \times \frac{\mathrm{d}u_{C}}{\mathrm{d}t}$. \
             On obtient~: \
             $u_{C}+RC\times\frac{\mathrm{d}u_{C}}{\mathrm{d}t}=E$

    \tcbitem (1pt)
    %----------------------------------------------------------------
    \tcbitem[raster multirow=4] 4.a)
    \tcbitem[raster multicolumn=2,raster multirow=2,blankest]
        \begin{tcbitemize}[raster rows=2,raster columns=2,
               raster height=\tcbtextheight,
               raster column 1/.style={width=\cwii},
               raster column 2/.style={width=\cwiii,fontupper=\bfseries}
            ]
            %--------------------------------------------------------
            \tcbitem On dérive $u_{C}$~: \
                     $\frac{\mathrm{d}u_{C}}{\mathrm{d}t}=
                     \frac{A}{\tau}e^{\frac{t}{\tau}}$.
                     On remplace dans l'équation différentielle~:
                     \begin{align*}
                         u_{C}+RC\times\frac{\mathrm{d} u_{C}}{\mathrm{d} t}=E
                         &\quad\Rightarrow\quad
                         A-Ae^{-\frac{t}{\tau}}+
                         RC\times\frac{A}{\tau}e^{-\frac{t}{\tau}}=E \\
                         &\quad\Rightarrow\quad
                         A+\left(\frac{RC}{\tau}-1\right)\times
                         Ae^{-\frac{t}{\tau}}=E
                     \end{align*}

                     L'équation doit être satisfaite quelque soit la valeur de
                     $t$, ceci implique~: $A=E$ 
            \tcbitem (0,5pt)
            %--------------------------------------------------------
            \tcbitem et \ $\left(\frac{RC}{\tau}-1\right)=0$
                     $\quad\Rightarrow\quad$ $\tau=RC$
            \tcbitem (0,5pt)
        \end{tcbitemize}
    %----------------------------------------------------------------
    \tcbitem[raster multirow=4] 4.b)
    \tcbitem[raster multicolumn=2,raster multirow=2,blankest]
        \begin{tcbitemize}[raster rows=2,raster columns=2,
               raster height=\tcbtextheight,
               raster column 1/.style={width=\cwii},
               raster column 2/.style={width=\cwiii,fontupper=\bfseries}
            ]
            \tcbitem Graphiquement~: $E=\qty{6}{\V}$
            \tcbitem (0,5pt)
            %--------------------------------------------------------
            \tcbitem Graphiquement~: $\tau=\qty{0.3}{\ms}=\qty{3e-4}{\s}$
            \tcbitem (0,5pt)
        \end{tcbitemize}
    %----------------------------------------------------------------
    \tcbitem 4.c)
    \tcbitem $\tau=RC$ $\quad\Rightarrow\quad$
             $C=\frac{\tau}{R}=\frac{\num{3e-4}}{100}=\qty{3e-6}{\F}=
             \qty{3}{\uF}$
    \tcbitem (1pt)
    %----------------------------------------------------------------
    \tcbitem 5)
    \tcbitem A
    \tcbitem (1pt)
\end{tcbitemize}

